\section{Theory}
\label{sec:theory}

The theory addresses two questions raised by the model definition. First, when
the observation is generated by a shared entrywise nonlinearity, can a
polynomial link approximate it without increasing the Tucker backbone rank?
Second, what multilinear-rank structure is induced by the tied powers
\(\cS^{\odot p}\)? We answer these as expressivity questions about the model
class.

\subsection{Approximation Under Nonlinear Observations}

\begin{theorem}[Polynomial generation]
\label{thm:poly-generation}
Suppose \(\cY = g(\cS^\star)\), where \(\cS^\star\) has Tucker rank at most
\((R_1,\ldots,R_N)\) and \(g\) is a polynomial of degree at most \(P\). Then
degree-\(P\) NTD-PL with the same Tucker rank can represent \(\cY\) exactly.
\end{theorem}

\begin{theorem}[Analytic link approximation]
\label{thm:analytic-approx}
Suppose \(\cY=g(\cS^\star)\), where \(\cS^\star\) has Tucker rank at most
\((R_1,\ldots,R_N)\) and \(\norm{\cS^\star}_\infty\le B\). If \(g\) is
analytic on the complex disk \(D(0,\rho)\) with \(B<\rho\), and
\(M_\rho=\sup_{|z|=\rho}|g(z)|\), then there exists a degree-\(P\) NTD-PL
model with the same Tucker backbone satisfying
\begin{equation}
  \frac{1}{\sqrt{M}}\norm{\cY-\widehat{\cY}}_F
  \le
  M_\rho
  \frac{(B/\rho)^{P+1}}{1-B/\rho},
  \qquad
  M=\prod_{n=1}^{N} I_n.
\end{equation}
\end{theorem}

Thus a finite polynomial link gives a direct approximation route for smooth
entrywise responses on a bounded latent range. The statement uses the simple
Taylor truncation bound; sharper polynomial approximation bounds can replace
it under stronger assumptions.

\subsection{Rank Expansion and Separation}

\begin{proposition}[Shared-backbone rank expansion]
\label{prop:rank-expansion}
Let
\(\cS=\tucker{\cX}{\bA^{(1)},\ldots,\bA^{(N)}}\)
have Tucker rank at most \((R_1,\ldots,R_N)\). For every integer \(p\ge 1\),
\(\cS^{\odot p}\) admits a Tucker representation with mode ranks at most
\begin{equation}
  \left(
  \binom{R_1+p-1}{p},\ldots,\binom{R_N+p-1}{p}
  \right).
\end{equation}
More explicitly, the ordinary Tucker factors of \(\cS^{\odot p}\) can be
chosen as Veronese lifts of the original factors. For
\(\bA\in\R^{I\times R}\), define
\(\bA^{\langle p\rangle}\in\R^{I\times D}\),
\(D=\binom{R+p-1}{p}\), by
\begin{equation}
  A^{\langle p\rangle}_{i,\alpha}
  =
  \prod_{r=1}^{R} A_{i,r}^{\alpha_r},
\end{equation}
where \(\alpha\in\mathbb N^R\) ranges over \(|\alpha|=p\).
Then \(\cS^{\odot p}\) factors through
\(\bA^{(1)\langle p\rangle},\ldots,\bA^{(N)\langle p\rangle}\).
\end{proposition}

This proposition makes the tying in
Section~\ref{sec:shared-backbone-interactions} explicit in standard Tucker
language: a degree-\(p\) channel is an ordinary Tucker tensor whose mode
factors are polynomial feature lifts of the same backbone factors. The link
therefore changes rank structure through tied features, not through an
independent high-order parameter block.

\begin{theorem}[Tensorial separation from standard Tucker]
\label{thm:tucker-separation}
Let \(N\ge 3\), \(R\ge 2\), \(p\ge 2\), and suppose \(I_n\ge D\) for every
\(n\), where \(D=\binom{R+p-1}{p}\). Then there exists an order-\(N\) tensor
\(\cY\in\R^{I_1\times\cdots\times I_N}\) that is representable by a
degree-\(p\) NTD-PL model with Tucker backbone rank \((R,\ldots,R)\), but
whose exact standard Tucker rank is
\begin{equation}
  (D,\ldots,D).
\end{equation}
Equivalently, a degree-\(p\) polynomial link can generate tensors whose
multilinear ranks are
\((\binom{R+p-1}{p},\ldots,\binom{R+p-1}{p})\) from a shared Tucker backbone
of rank \((R,\ldots,R)\).
\end{theorem}

The proof in Appendix~\ref{app:proofs} uses a CP-diagonal Tucker backbone and
generic full-rank Veronese feature matrices. The theorem is an existence
separation, not a universality claim over all high-rank Tucker tensors: it
identifies a structured high-rank subset generated by a low-rank shared
backbone.

Together, the results clarify the scope of NTD-PL. Polynomial links can
approximate shared nonlinear observation maps while keeping the same latent
Tucker rank, and their powers induce Veronese-lifted Tucker factors. This
explains how a scalar link can change multilinear rank structure without
introducing untied high-order tensor parameters.
